\section{Experiments}

All development ran against the development set of BIRD~\cite{bird}, 1,534
questions that the benchmark splits into 925 simple, 464 moderate, and 145
challenging. The test set is private, and no part of the development loop
touched it. Quality is reported as execution accuracy, scored exactly as every
gate decision was scored (Section~\ref{sec:methodology}): a prediction counts
as correct when the rows it returns, compared as a set, agree with the rows of
the reference query.

Three generating models appear in the measurements. Claude Opus~4.8 is the
model the pipeline was developed against, and every acceptance and rollback
decision was measured on it. Claude Opus~4.6 is an earlier version of the same
model. GigaCode~1.6 comes from a different model family and a different
developer. Each model appears in two configurations. Single-shot issues
one call carrying the instructions, the database schema, and the question, and
takes the returned query as the answer. Agent-built runs the 17-step pipeline
that the quality gate admitted. Table~\ref{tab:results} reports the six
configurations.

\begin{table}[t]
\caption{Execution accuracy (EX, \%) on BIRD by difficulty tier. Single-shot is
one call to the model, carrying the instructions, the database schema, and the
question. Agent-built is the 17-step pipeline the quality gate produced, run
unchanged on all three models. A dash marks a score that was not
measured.}
\label{tab:results}
\centering
\footnotesize
\setlength{\tabcolsep}{4pt}
\begin{tabular}{@{}lrrrrr@{}}
\toprule
Configuration & Simp. & Mod. & Chall. & Dev & Test \\
\midrule
\multicolumn{6}{@{}l}{Claude Opus 4.6} \\
\quad single-shot & --- & --- & --- & 68.8\rlap{$^*$} & 70.2\rlap{$^*$} \\
\quad agent-built & 77.1 & 64.7 & 62.8 & 72.0 & --- \\
\addlinespace[2pt]
\multicolumn{6}{@{}l}{Claude Opus 4.8} \\
\quad single-shot & 75.0 & 64.7 & 64.8 & 70.9 & --- \\
\quad agent-built & 80.7 & 72.0 & 61.4 & 76.2 & 81.3\rlap{$^\dagger$} \\
\addlinespace[2pt]
\multicolumn{6}{@{}l}{GigaCode 1.6} \\
\quad single-shot & 66.2 & 52.2 & 43.4 & 59.8 & --- \\
\quad agent-built & 73.5 & 56.9 & 52.4 & 66.5 & --- \\
\bottomrule
\end{tabular}

\smallskip
\begin{minipage}{\columnwidth}
\footnotesize\raggedright
$^*$\,Official BIRD leaderboard baseline for Claude Opus~4.6, quoted
from~\cite{bird-leaderboard}.\\
$^\dagger$\,Our submission, scored by the BIRD team on the private test set
and listed on the leaderboard~\cite{bird-leaderboard}.
\end{minipage}
\end{table}

The comparison that answers RQ1 is the same-model pair. One call to Claude
Opus~4.8 scores 70.9 on the development set, and the pipeline of steps admitted
by the gate, built around that same model, scores 76.2. The model is identical
in both rows, so the 5.3 points come from the context the pipeline builds for
the model and from the repairs it applies to the query the model returns. Seventeen steps written
by an agent and admitted on measured effect alone compose into an artifact that
is better than the model called directly.

RQ3 asks whether that quality belongs to the model it was measured against.
Re-run with no step changed, the pipeline lifts Claude Opus~4.6 from 68.8 to
72.0, a gain of 3.2 points, and GigaCode~1.6 from 59.8 to 66.5, a gain of 6.7.
The gain reproduces on all three models and spans 3.2 to 6.7 points, and it
does not track the strength of the base model: the largest gain belongs to the
weakest of the three, and the cross-vendor model gains more than the earlier
version of the model the steps were tuned on. Both transfer measurements are
development-set measurements, since the test set was submitted only for the
model the pipeline was developed against. There the agent-built pipeline scores
81.3---third on the BIRD leaderboard as of August 2026 and 0.6 points off
first place~\cite{bird-leaderboard}---which is external evidence that the
accepted artifact holds up on questions that no development decision could have
been fitted to. The single-shot baseline also rises between the two sets, by 1.4
points on Claude Opus~4.6, so part of the difference between 76.2 and 81.3
belongs to the question sets rather than to the pipeline.

RQ2 asks which defects the mining surfaced and what each repair is worth. The
answer is read off ordinary runs. Because every step records what it received
and what it returned, the run that scores 76.2 can be decomposed after the fact
without a single ablation rerun. From the 51,706 recorded step spans of that
run, each question's query was reconstructed at every stage of the chain,
executed against its database, and compared with the reference by the same
comparator the run itself used, 3,013 distinct executions in total. A step
\emph{saved} a question when the query entering it was wrong and the query
leaving it was right, and \emph{broke} the question in the reverse case. The
effect is local to the step's position in the chain, so one question can be
saved by one step and broken by another, which happened four times. As a
consistency check, the recomputed final verdicts agree with the verdicts
recorded by the run on all 1,532 questions the reconstruction could
re-execute.

\begin{table}[t]
\caption{Contribution of each step that rewrites the query, reconstructed from
the recorded inputs and outputs of one ordinary run: Claude Opus~4.8 with the
agent-built pipeline, the 76.2\% row of Table~\ref{tab:results}. Fired counts
the questions whose query the step rewrote, saved and broke count the rewrites
that turned a wrong answer right and a right answer wrong, and net is their
difference. Step~9 sets the baseline of 1,137 correct answers out of 1,534.
Steps~1--8 do not touch the query and admit no per-question attribution.}
\label{tab:contrib}
\centering
\footnotesize
\setlength{\tabcolsep}{3pt}
\begin{tabular}{@{}rlrrrr@{}}
\toprule
\# & Step & Fired & Saved & Broke & Net \\
\midrule
9 & \texttt{generate\_sql} & 1,534 & --- & --- & --- \\
10 & \texttt{judge\_sql} & 46 & 13 & 6 & $+7$ \\
11 & \texttt{fix\_join\_type\_mismatch} & 0 & 0 & 0 & 0 \\
12 & \texttt{add\_null\_filters} & 0 & 0 & 0 & 0 \\
13 & \texttt{audit\_select\_list} & 11 & 3 & 3 & 0 \\
14 & \texttt{repair\_empty\_result} & 6 & 4 & 0 & $+4$ \\
15 & \texttt{strip\_distinct\_count\_id} & 51 & 25 & 8 & $+17$ \\
16 & \texttt{uppercase\_yes\_no} & 3 & 1 & 0 & $+1$ \\
17 & \texttt{cast\_real\_division} & 27 & 4 & 1 & $+3$ \\
\midrule
\multicolumn{5}{@{}l}{Net after generation} & $+32$ \\
\bottomrule
\end{tabular}
\end{table}

Table~\ref{tab:contrib} reports the steps that rewrite the query. The eight
steps before generation prepare the generator's input and never edit a query, so the
table omits them and their activity is reported here: table selection, schema
building, column resolution, and schema focusing run on all 1,534 questions,
cell-value hints are found on 867, conflicting-concept annotations are injected
into 1,201, the reversed linker adds tables on 16, and the table auditor
expands a selection once. Of the total gain of 5.3 points over single-shot
generation, 3.2 are already present when the generator returns its query, which
is correct on 1,137 questions or 74.1\%, and the remaining 2.1 points are the
net $+32$ answers added by the steps that check and repair that query. More
than half of those 32 come from one step---the removal of a \texttt{DISTINCT}
inside a count over a key column---which is a database-specific counting
convention rather than a reasoning failure. The judge opened its candidate
pool on 118 questions and rewrote 46 of them. It nets $+7$ but works noisily:
only 19 of the 46 rewrites changed a verdict, 13 for the better and 6 for the
worse, and 22 replaced one wrong query with another. The two steps that never fire found nothing to repair
across all 1,534 questions, because the failures that motivated them are
prevented earlier in the chain in this version of the pipeline.

Sliced by difficulty, the gain is uneven. Against the
single-shot baseline on Claude Opus~4.8, the pipeline adds 5.6 points on simple
questions and 7.3 on moderate ones, and loses 3.4 on challenging ones, from
64.8 to 61.4. Across all Claude Opus configurations in
Table~\ref{tab:results} the challenging tier stays between 61 and 65, whether
the query comes from one call or from the pipeline. We read all 56 challenging
misses of the final run by hand. In most of them the prediction is a
defensible reading of the question that differs from the reference only in
which columns it projects or in which quantity it counts, in several the
reference contradicts the evidence hint attached to the question, and a
minority are outright reasoning errors. No rewrite aimed at any of these
classes survived measurement. What the gate harvests, on this model, are
convention-level defects that recur across many questions, and those are
concentrated in the simple and moderate tiers. The challenging tier is not
inert everywhere. On GigaCode~1.6 the pipeline adds 9.0 points there, from 43.4
to 52.4, more than it adds on simple questions (7.4) or moderate ones (4.7). A
tier that does not move for a strong model moves for a weaker one, which places
the ceiling with the model rather than with the pipeline.
